\documentclass[11pt, a4paper]{article}

\usepackage[margin=1in]{geometry}
\usepackage{amsmath}
\usepackage{amssymb}
\usepackage{xcolor}
\usepackage{enumitem}
\usepackage{url}
\usepackage{bm}

% --- Custom commands and environments ---
% Use this command for reviewer headings
\newcommand{\reviewer}[1]{\subsection*{#1}}

% Use this environment for your answers.
% You can easily change the style of all your answers by modifying this definition.
% For example, to remove the color, just delete the \color{...} line.
\newenvironment{answer}{%
      \begin{quote}%
            \color{blue!80!black}%
            \textbf{Answer:}%
            % Add your response here.
            }{%
      \end{quote}%
}

% Use this environment for your private notes.
% To hide all notes, comment the definition below and uncomment the alternative.
\newenvironment{notes}{%
      \begin{quote}%
            \color{gray}%
            \textbf{Notes:}%
            % Add your notes here.
            }{%
      \end{quote}%
}

% Mark changes/additions by the AI assistant
\newcommand{\chadded}[1]{\textcolor{blue!70!black}{\textbf{(chadded:}~#1\textbf{)}}}

% Alternative definition to hide all notes:
% \newsavebox{\hidenotes}
% \newenvironment{notes}
%   {\begin{lrbox}{\hidenotes}}
%   {\end{lrbox}}
% --- End of custom commands ---

\begin{document}
\title{Efficient Matrix-Free Implementation of Ghost Penalty Stabilization in Cut Finite Element Methods}

\section*{Response to the reviewers}
I would like to thank the editor and the reviewers for their constructive comments and suggestions, which have helped
me to improve the manuscript. I have revised the manuscript accordingly and provide a point-by-point response to the
reviewers' comments below.

I am happy to say that the most of the presented methods  are already included in 9.7 release of deal.II, the paper is
now cited in the manuscript.

I removed the section containing convergence results from the manuscript. Figure~1 was corrected to include one
previously missing ghost‑penalty face.
Previous versions of the manuscript contained a more detailed discussion of the convergence properties of the proposed
method, this has been removed as it was redundant with prior work.
The styles of all plots have been improved for clarity and consistency.
Two plots were combined into one figure with subfigures (Figure 6 in the current version of the manuscript).
Bibliographic entries were fixed and new citations to the Shifted Boundary Method were added to better situate the
present work within related approaches.

\section*{Editor's and Reviewers' comments}

\reviewer{Reviewer \#1}

This paper proposes an efficient matrix-free implementation of ghost penalty stabilization within the Cut Finite
Element Method. The author explores the tensor-product structure of basis functions, which enables the simplification
of evaluating ghost penalties to 1D matrix-vector products. The implementation targets high-order finite element
methods and is based on the deal.II library. I recommend publishing the article after incorporating some minor updates
(see the comments below) to enhance the presentation.

\begin{enumerate}[label=\textbullet]
      \item Abstract would benefit from providing more details on the necessity of the presented work.
            \begin{answer}
                  Added a sentence in the abstract.
            \end{answer}

      \item "For simplicity of presentation, we consider a 3D case with a hexahedral element; however, the reasoning
            can be easily transferred to the 2D case with quadrilateral elements." Isn't it for generality rather than
            simplicity?
            \begin{answer}
                  The idea is that it avoids writing products over dimensions, and the 3D is a specific case and not
                  fully general. Hence, I believe simplicity is the right word here.
            \end{answer}

      \item Section 2.4: The derivation relies on the fact that the cells are aligned with the coordinate axes.
            Please, emphasise it in the text and the abstract.
            \begin{answer}
                  Done.
            \end{answer}

      \item Examples: Are the cut-integration rules precomputed, or are they computed on the fly? Is it possible to
            compare the time required for precomputing them versus computing on the fly?
            \begin{answer}
                  I have clarified that the
                  cut-integration rules are precomputed and stored, and that this breaks the tensor product
                  structure.

                  The on-the-fly computation of these rules is computationally expensive, and the underlying algorithms
                  often involve significant branching, which is not compatible with SIMD. Since the amount of data
                  produced (quadrature point locations and weights) is manageable, pre-computing and storing these
                  rules is a more efficient strategy.
            \end{answer}
\end{enumerate}

\reviewer{Reviewer \#2}

The manuscript presents an efficient implementation of ghost-penalty stabilization for CutFEM discretizations. It
relies on matrix-free operator evaluation for the physical PDE terms, Nitsche's method for homoegeneous Dirichlet
boundary conditions, and a Kronecker form of ghost penalty stabilization specific to Cartesian mesh elements. The
manuscript describes the underlying mathematical concepts, explains an efficient implement, presents a test of
verification and evaluates the performance on various scenarios in 2D and 3D, with the main focus on 3D.

The manuscript covers an important topic, because unfitted finite element methods have been known to offer many
advantages for certain applications, such as simplified meshing in context of complex geometries or moving domains, but
the efficient implementation has been lagging, especially for higher orders and in 3D. The chosen ingredients are
state-of-the-art and based on the widely used deal.II finite element library, providing a credible software
environment. The obtained results are also very encouraging. At the same time, many of the concepts, software
ingredients and evaluation steps have been presented in a previous publication by Bergbauer et al. (arXiv preprint from
April 2024, published online in SISC in May 2025, \url{https://doi.org/10.1137/24M1653689} ). Given the similar scope
and timing, I suggest to clearly elaborate on the novelty of the present work. I believe that the Kronecker expression
of the ghost penalty term is worth to be published and evaluated, but I also believe that the manuscript should
concentrate on that part primarily, rather than repeat what has been said before. Furthermore, the manuscript is not
prepared very carefully, making several parts losely connected, jumping around between different lines of reasoning,
containing many avoidable errors and similar deficiencies that make reading very difficult. Below, I provide a long
list of comments, some bigger in scope, some just minor remarks, that should be addressed before I can consider
recommending the manuscript for publication. Also I suggest to carefully read everything again, because I could have
missed errors or imprecise statements. But the main goal of a revision should be to:
\begin{itemize}
      \item[a)] make sure to point out the novelty against the contribution [3], and
      \item[b)] perform experiments to assess the performance of the proposed method side-by-side against the published
            software implementation of [3], to evaluate the novel parts of the contribution.
\end{itemize}

\begin{answer}
      Thank you for the detailed feedback and the opportunity to clarify the contributions of our work. I have revised
      the manuscript to more clearly distinguish our contributions from those in [3].
\end{answer}

Ideally, the benchmark and proposed code could be made open-source, or possibly even integrated into the deal.II finite
element library. This could further contribute to verification and software sustainability, but I do not see it as an
absolute requirement.
\begin{answer}
      Parts of the code has already been integrated into deal.II, specifically assembly of the 1D ghost penalty
      stabilization
      term and numbering of DoFs on two adjacent cells. See PR18413 (numbering)
      PR18361 (1D matrices and Kronecker product for classical assembly), PR18061 (additional ghost cell in
      matrix-free).
      The efficient implementation of Kronecker product is available in deal.II for a long time as a part of
      matrix-free framework.

      There is also a PR18746 that updated the step-85 tutorial, however it was decided that the new way of applying
      ghost penalty stabilization should be presented as a possible extension to keep the simplicity. The results
      presented in this PR confirm the correctness of the implementation.
\end{answer}

\bigskip
\noindent\textbf{List of comments:}

\begin{enumerate}[label=\textbullet, wide, labelindent=0pt]
      \item p1 first paragraph: "when elements are small cuts" and "very small subdomains" not phrased accurately.
            \begin{answer}
                  Clarified the phrasing.
            \end{answer}

      \item p1, 2nd paragraph: The statement "The implementation of the ghost penalty is somewhat vexing because it
            involves evaluation of high-order derivatives." is not entirely correct because high-order derivatives are
            not the only
            possible choice. One can use volume-based ghost penalty methods (later referenced on page 2) that are based
            on the L2
            projection on patches, which was to the best of my knowledge initially described in Sec. 4 of Burman, 2010
            (Ref [4] in
            the manuscript), see also later uses in J. Preuss, 2018 (Ref [33]) or Bergbauer et al., SISC, 2025 (Ref
                  [3]). Regarding
            the text later on page 2, paragraph 4, where it is stated that macro-patch stabilization may lead to
            locking: The end
            result of that stabilization approach is roughly equivalent to the face-based derivatives in that
            polynomial spaces of
            adjacent elements get linked, so the locking part should apply to most ghost penalty approaches, including
            the one in
            the manuscript, in my opinion. In fact, Ref [2] that discusses the issue of locking attributes a more
            strict limit to
            face-based ghost penaly than to the volume-based (called bulk-GP in that work), see their Sec 4.3, and Sec.
            2.5 of Ref
                  [8] (Burman at al.). The crucial part is not bulk or face derivative coupling, but the definition of
            patches. Note that
            the selected approach visualized in Fig 1 chooses coupling between all faces, so it is susceptible to
            locking in the
            sense of [2, 8]. This topic should be clarified and more appropriate context given. Besides the erroneous
            classification of face-based and bulk-based stabilization with respect to locking, it is probably also not
            crucial to
            be locking-free for the present work because a fixed parameter is used, rather than going towards infinity.
            \begin{answer}
                  We now acknowledge volume-based alternatives, clarify that face-based methods (including our own) can
                  also be susceptible to locking, and note that the patch definition is the crucial factor, citing the
                  recommended
                  literature.
            \end{answer}

      \item p1, 2nd paragraph: This paragraph discusses how to express the ghost penalty operator as a Kronecker
            product (which, to the best of my knowledge, has not been done before), but overall the efficient
            evaluation
            was
            investigated in detail in a previous publication [3]. While there is a small remark comparing to [3] on
            page
            2,
            paragraph 2, in my opinion the present work should be contrasted against that contribution, identify
            analogies and
            differences to give the reader a clear picture of the respective metrics, see also the general comment
            above.
            \begin{answer}
                  The paragraph has been expanded to better contrast with [3].
            \end{answer}

      \item p2, 2nd paragraph, 1st line: There are two references to ref [1].
            \begin{answer}
                  Corrected.
            \end{answer}

      \item p2, paragraph on TraceFEM: Is this crucial for the manuscript and is it used later on that I just missed?
            If not, I suggest to shorten this part as it might be confusing to the reader when talking about a
            different
            topic.
            \begin{answer}
                  Paragraph shortened to a single sentence, moved at the end of paragraph starting with "Ghost penalty
                  methods..."
            \end{answer}

      \item p2, 4th paragraph: The text states "Larson et al. [28]" but there are only two authors in [28], Larson
            and Zahedi, where the common citation style would be "Larson and Zahedi [28]". In general, the citation
            style
            of
            literature results could be made more uniform throughout the manuscript. Let me mention one single
            additional
            problem,
            besides several more that could be listed: the citation of [35] mis-spells the first author name (Sch\"oder
            vs the
            spelling Schoeder in the reference list, the latter is also the one I find online), and does not use Stokes
            but
            acoustics according to the paper title.
            \begin{answer}
                  Done.  I have also corrected other citations. The bibtex compilation was failing and some entries
                  authors were replaced with ----, this is also fixed.
            \end{answer}

      \item p2, last paragraph: The text on matrix-free methods is more or less a repetition of the first paragraph
            on the same page. Those should be merged in a single place.
            \begin{answer}
                  Merged the two paragraphs into one.
            \end{answer}

      \item p4, definition of Nitsche penalty \(\gamma_D = 30 k (k+1)\): How is the factor 30 derived? In general,
            one would use a 1/h scaling here to ensure coercivity, see, e.g., Warburton and Hesthaven, "On the
            constants
            in
            hp-finite element trace inverse inequalities", CMAME 192, 2003 (doi:10.1016/S0045-7825(03)00294-9), and
            related work on
            trace estimates. Please clarify the choice and to which extent this is robust for any mesh size h.
            \begin{answer}
                  Done. The missing part was 1/h scaling to the Nitsche penalty term in the weak form and in
                  the algorithm description.

                  I was experimenting with some preconditioning strategies and 30 behaved
                  well for patch smoothing. It does not affect the results presented in the manuscript, but for the
                  sake of completeness I included it.

                  The paper by Warburton and Hesthaven is now cited.
            \end{answer}

      \item p5, Sec 2.4: The initial text mentions \(\xi\) as the reference coordinate, whereas the two first
            formulas in this text use the coordinate \(x_i\). This is confusing. I suggest to clarify the notation in
            terms of
            reference coordinates and physical coordinates and make sure to be consistent.
            \begin{answer}
                  Done. I have changed \(\xi\) to \(x_1\) in the said formula.
            \end{answer}

      \item p5, text describing Fig 2: This is a 2D figure, but the text around it talks about 3D being the case
            discussed. This should be clarified and made uniform.
            \begin{answer}
                  Done. I have added a sentence to the caption of Figure 2.
            \end{answer}

      \item p6, equation (1): The bilinear form \( ( , )_F \) represents an integral over a face F in the first and
            second line. However, the third line starts factoring out some expressions, but the integration domain
            still
            remains F
            on each constituent. While it is intuitively clear what happens, the notation does not quite match; this
            should be
            corrected, e.g. by stating the underlying integrals.
            \begin{answer}
                  Changed the intergration, use $F^1$ to denote the face in reference coordinates. Face $F$ was already
                  described as a product of 1D sets.
            \end{answer}

      \item p6, paragraph before formula (3): Isn't the size of the ghost penaly matrix (2k+1)x(2k+1), rather than
            (2k-1)? See also many other places in the text.
            \begin{answer}
                  Done. The size of the ghost penalty matrix is now stated as
                  $(2k+1) \times (2k+1)$.
            \end{answer}

      \item p7, formula after Eq. (4): Isn't the mass matrix scaled by h rather than \(h^{-1}\), because the factor
            arises from the change of size of the integration domain (0,1) in reference coordinates to an interval of
            length h in
            real coordinates? Please check and make sure to present the correct formula, see also the use of formula
            (5).
            \begin{answer} Corrected. Eq (5) was correct.
            \end{answer}

      \item p7, text immediately after (5): "We store the matrices precomputed mass matrix h M\textasciicircum1 and
            the penalty matrix" is not grammatically correct.
            \begin{answer}
                  Done. The sentence has been corrected.
            \end{answer}

      \item p8, 1st paragraph: The text "The second step applies the mass matrix M h in the direction orthogonal to
            the face, and the third step applies the mass matrix in the direction parallel to the face." does not
            appear
            correct,
            as both mass matrices arise from the tangential space on the face, whereas the ghost penalty matrix is the
            one aligned
            with the face.
            \begin{answer}
                  The description of the mass matrix application has been corrected
            \end{answer}

      \item p8, 1st paragraph, complexity: Since k is not very large in the examples shown later (with k=1,2,3
            evaluated; only Fig 4 has some general scaling to higher k), the constants in front of the big-O notation
            should be
            given as well, or other arguments in how the approach differs from [3], for example.
            \begin{answer}
                  Done. I have replaced the big O notation with the explicit formula for the computational
                  cost.
            \end{answer}

      \item p8, Sec 2.5.2: I do not really understand the 2nd part in the enumerated list, "Quadrature loop: Queue
            \(\nabla u_q\) for integration." - shouldn't this also involve the test function, Jacobian determinant and
            quadrature
            weight? I see a similar presentation in Algorithm 1. Also, it is not really clear how the summation of the
            gradient of
            test function with index i is treated in steps 2 or 3. I believe I see how it would get implemented in a
            code, but the
            presentation in the text should be improved. The same holds for formulas (12)-(14) and their relation to
            the
            text.
            Furthermore, in Algorithm 1, line 22, not just the queued gradients but also queued values should be
            integrated. Here,
            the imprecise notation gets particularly striking, because the part that gets queued is not the value or
            gradient on
            its own, but the content to be multiplied by the value or gradient of the test function!
            \begin{answer}
                  Jacobian determinant, quadrature weight, and test function are applied during the
                  integration
                  step.
                  While Jacobian determinant and quadrature weight are just scaling factors, multiplication by test
                  function is done via sum factorization.

                  The text has been clarified to explain the evaluation and integration steps.
            \end{answer}

      \item Regarding Secs 2.5-2.7 and the related parts of the manuscript: These are not new algorithms but follow
            previous works closely. Besides the references already given on matrix-free evaluation with sum
            factorization, it would
            be good to also list some classical textbook on spectral element methods, such as Deville, Fischer and Mund
            (2002) or
            Kopriva (2009). Furthermore, elaborate on what is new.
            \begin{answer}
                  Added citations to the recommended literature and clarified that the treatment of interior cells
                  follows the standard matrix-free approach.
            \end{answer}

      \item p8, bottom: The sentence "Note the the cost the cell evaluation grows in the same rate as the cost of
            evalution of the ghost penalty," appears to be either using the wrong punctuation or be incomplete, and it
            contains
            errors.
            \begin{answer}
                  Corrected.
            \end{answer}

      \item p9, itemized list at top of page: Here, the symbol J is mentioned to denote the Jacobian matrices, but
            earlier formulas use J for the determinant of the Jacobian. Use a consistent notation. Also, for formula
            (16)
            the
            common notation would be to use \(J^{-T}\) (inverse Jacobian, transposed), but of course this depends on
            how
            exactly
            the rows and columns of the Jacobian matrix are defined. This should be clarified as well. (Note that the
            mesh cells
            are Cartesian, so the implementation is correct also without transposition, but it is not good style to
            confuse
            readers.)
            \begin{answer}
                  Changed the notation to bold $\bm{J}$ for the Jacobian matrix and non-bold $J$ for its determinant.
                  Changed the formula to use \(\bm{J}^{-T}\).
            \end{answer}

      \item p9, numbered list on the lower part of the page: I am not convinced by the order of the steps. It seems
            as if steps 1 and 3 belong together and steps 2 and 4 belong together, which makes it unclear as to why
            they
            are not
            ordered appropriately. In any case, these are also not three steps as claimed at the top of the list. This
            should be
            clarified or fixed. In any case, I believe all of Sec 2.6 is done in the same way as in Ref [3].
            \begin{answer}
                  From mathematical point of view that is correct, but from implementation point of view I do it this
                  way because this allows us to read DoFs of a cell only once and then directly use them to compute all
                  quantities needed
                  for the cell and associated immersed boundary. The memory used for fetching DoFs is used then to
                  accumulate the results of integration.

                  This has been clarified in the text.
            \end{answer}

      \item p9: Regarding the integration on cut cells / faces, it should be clarified whether some tensor product
            structure is used, and if this affects the complexity estimates I mentioned above. Note that the text at
            the
            bottom of
            page 10 indicates a higher complexity, but this should be made more explicit at this place.
            \begin{answer}
                  Cut cell integration is not tensor product, this is now clarified in the text.
            \end{answer}

      \item p10, 2nd paragraph: The text "As we expect this [the matrix-vector product] to be the most
            computationally intensive part of any matrix-free solver for CutFEM, we aim to utilize the hardware as
            efficiently as
            possible." might be misleading. In a matrix-free context, the most challenging part is usually to find
            efficient
            preconditioners, in particular to find preconditioners that have similar cost as the (excellent) algorithms
            derived by
            present contribution.
            \begin{answer}
                  - Rephrased the sentence to be more precise.

                  - Added citation to patch smoothing.
            \end{answer}

      \item p10, 5th paragraph: The text "The impact of this imbalance is limited as the number of faces with ghost
            penalty is small compared to the total number of cells" is misleading without context, because in 3D the
            number of
            faces is not that low (think about the 25 sphere case), and what matters is the computational cost of the
            face penalty
            matrices compared to the other ingredients in the algorithm and the memory access.
            \begin{answer}
                  Rephrased.
            \end{answer}

      \item p10, 6th paragraph: This starts with "Finally the ghost penalty operator is applied to the faces." but
            this is already what gets discussed in the previous paragraph. The text should be made concise, group one
            line of
            argument into the same paragraph instead of jumping back and forth, and avoid repetition.
            \begin{answer}
                  Done, the order of the paragraphs has been changed.
                  The repetition has been removed.
            \end{answer}

      \item p10, last paragraph: The description of load balancing belongs to the MPI parallelization in my opinion
            (4th paragraph). This is confusing.
            \begin{answer}
                  Done. Text in that entire subsection was be restructured.
            \end{answer}

      \item p12, 1st paragraph: Isn't the compile flag "-O3" rather than "O3"?
            \begin{answer}
                  Corrected.
            \end{answer}

      \item p12, 2nd paragraph: "interior of unit sphere" - does this mean the unit ball, or is it a surface problem
            (2D surface embedded into a 3D space) that gets solved here? In the latter case, how are the details?
            \begin{answer}
                  I have changed "interior of the unit sphere" to "unit ball".
            \end{answer}

      \item p12, Sec 3.1: This convergence test is very similar to what many works on CutFEM have done before.
            \begin{answer}
                  I agree that section was indeed redundant and has been removed.

                  The goal was to illustrate the correctness of the implementation, but since the kronecker evaluation
                  of the ghost penalty operator will become a test inside deal.II (Implementation form PR18746 will be
                  moved to tests) this is not needed.
            \end{answer}

      \item p13: Ref [3] appears to have the experimental code available as open-source, so a comparison of the
            actual code in a side-by-side experiment should be made (in the SISC publication of [3] I found the
            reference
            \url{https://github.com/bergbauer/fe_point_evaluation_benchmarks} ). Furthermore, it is crucial to explain
            how the
            experiments shown in this section are similar or differ from the benchmarks in earlier work. As they are
            done
            right
            now, several pieces of information feel a bit repetitive.
            \begin{answer}
                  I was not aware of this repo.

                  However in this repo ghost penalty is not implemented.
                  FEEvaluation and  FEPointEvaluation are exactly the classes used in [3] so the comparison was already
                  done, just the description was inaccurate.

                  The micro-benchmarks in the paper are done in a way that emphasizes the number of floating point
                  operations and the memory accesses should not interfere.

                  The benchmarked code is the same. I just add ghost penalty on top of it.
            \end{answer}
      \item p13: The text first introduces unknown names "FEEvaluation" or "FEPointEvaluation", and some
            abbreviations "FEEval" and "FEPointEval" in the figures; later, it is explained that those come from
            deal.II.
            The text
            should be made more clear to not have the reader guess first where the name come from, go around in the
            previous
            sections of the text, only to find out a few sentences later where they come from. (It might be just me
            reading the
            text piece by piece and wondering before I go on with reading, but this could be seen as a remark to work
            on
            the
            organization of the information.)
            \begin{answer}
                  Done. I have clarified the text and used \texttt{texttt} for the class names.
            \end{answer}

      \item p13, experiment to illustrate relative timings: I do not understand the experiment in its entirety. The
            "reference" time lists 0.077 us, which is less than a single miss to main memory on modern CPUs. There is
            some vague
            statements above performing SIMD vectorized work with a magic number of "8" that is used as a division
            factor
            (albeit
            the listed CPU of AMD appears to not support 512-bit vectorization; maybe this is a single precision
            experiment or an
            experiment with many elements?). Furthermore, I am puzzled whether typical timers on CPUs can provide
            sub-microsecond
            resolution, because the timing would interfere with out-of-order execution and similar features within the
            CPU. Is this
            experiment repeated several times to get reliable numbers? Are there statistical deviations to take into
            account? It
            might be good to read the text on "Scientific Benchmarking of Parallel Computing Systems" by Hoefler and
            Belli, 2015,
            \url{https://doi.org/10.1145/2807591.2807644} , for the state-of-the-art concerning timing results.
            \begin{answer}
                  I fully agree with this comment, the experiment was planned to avoid those issues, however the
                  detailed
                  description was not provided in the paper.

                  The experiment is indeed repeated many
                  times to get reliable numbers, the total time is measured over all repetitions, the reported values
                  are averages. I have clarified the text regarding the experimental setup.

                  The experiment is done in double precision and indeed vectorization is done with AVX2 (256 bit).

                  I looked into the referenced text, it is indeed useful, I added a citation.

                  This is now corrected in the text and the plots are adjusted accordingly.
            \end{answer}

      \item p13/14: It seems as if the experiment of Fig 4 (relative timings) is a sequential one, whereas the
            matrix-vector multiplications (vmult) of Figs 5 and 6 are done in parallel. The experimental setup of the
            tests should
            be made more clear, the generic text at the beginning of Sec 3 does get blurried by the experiment of Fig
            4.
            It might
            make sense to split up that part into a different subsection.
            \begin{answer}
                  This benchmark is indeed in a sequential setting, in contrast to the parallel benchmarks in Figures
                  5 and 6. The purpose of this specific experiment is to measure the performance of the individual
                  components of the
                  matrix-vector product in isolation, making it a compute-bound micro-benchmark. The goal of Figure 4
                  is to illustrate
                  the computational complexity of the different components and to verify the theoretical complexity
                  estimates discussed
                  earlier.
                  It is now clarified in the text.
            \end{answer}

      \item p13: The metrics of DoFs/second is not clearly defined. The text on page 13 talks about "average time per
            degree of freedom", but the figures show the reciprocal (DoFs/s). Also, it is not clear whether DoFs
            outside
            the domain
            but part of the fictious domain are counted as well, or whether those parts do not carry any degrees of
            freedom. Does
            the code handle this setting in an efficient way? This should be clarify, especially for the 25-ball case
            later in the
            subsection.
            \begin{answer}
                  The metric is indeed DoFs/s, I have corrected the text accordingly.

                  The DoFs outside the domain are not counted, the text has been clarified to explain this.
            \end{answer}

      \item p14, 2nd paragraph: "For a fixed 3D mesh consisting of obtained by refining a single cell 5 times" - this
            is strange grammatically and also from the actual mesh I do not get the point. Is the "5" the number of
            uniform
            refinement steps, doubling the number of elements per direction in each step, or is it adaptive refinement?
            If adaptive
            meshes get used, please explain are they treated with ghost penalties and the resulting differences in mesh
            sizes (and
            hanging nodes).
            \begin{answer}
                  Done. I have added a sentence to clarify the timing breakdown.
            \end{answer}

      \item p15/16: The timing breakdown for the experiment (Fig 9) should be made more clear: Does the "MPI" time
            include possible load imbalances? How are the timings distributed over the 32 MPI ranks (and is "32" the
            right number
            for the experiment)? Is the timing recorded for the cut cells, which is shown to dominate, explained by the
            low-level
            experiment of Fig 4? And in turn, can the recorded timings be explained by the computational complexities,
            or
            are there
            hidden factors? More precise numbers and relations between the experiments should be given to make the
            results
            understandable for the reader and create an added value compared to previous works.
            \begin{answer}
                  - MPI time includes load imbalance, yes.

                  - 32 is the number of MPI ranks for the experiment.

                  - The timings are summed up over all MPI ranks, so MPI time includes the time wasted due to load
                  imbalance.

                  - The timing recorded for the cut cells is indeed explained by the low-level experiment of Fig 4

                  The text has been clarified to explain those points and relate the experiments from Fig 4 and Fig 8.
            \end{answer}

      \item p16, conclusions: The text compares against sparse matrices, which is ok, but it does not mention that
            previous works have obtained similar complexity/cost benefits (Ref [3]), and that convergence
            rates/flexibility of
            CutFEM as a concept are simply reproducing the work by Burman, Larson and many co-workers.
            \begin{answer}
                  Conclusions was rewritten to clarify the novelty of the present work.
            \end{answer}

      \item General appearance: Some equations are numbered in tight sequence (e.g., (7)-(9)), whereas others are
            not. I suggest to either number all equations, or to number only those that the text refers to.
            \begin{answer}
                  All equations are now numbered.
            \end{answer}
\end{enumerate}

\end{document}